\documentclass[10pt,landscape]{article}
\usepackage[utf8]{inputenc}
\usepackage[T1]{fontenc}
\usepackage{amsmath, amssymb, amsthm}
\usepackage{multicol}
\usepackage{geometry}
\usepackage{hyperref}
\usepackage{xcolor}
\usepackage{enumitem}

% Set page geometry for a landscape 3-column layout
\geometry{top=0.5in, bottom=0.5in, left=0.5in, right=0.5in}

% Custom settings for columns and lists
\setlength{\columnsep}{0.4in}
\setlength{\columnseprule}{0.4pt}
\setlist[itemize]{noitemsep, topsep=0pt, leftmargin=*}
\setlist[enumerate]{noitemsep, topsep=0pt, leftmargin=*}

\title{Cheat Sheet}
\author{}
\date{}

\begin{document}

\begin{multicols*}{2}

\section*{0. Math and Probability}

\textbf{Exponentials}
\begin{itemize}
    \item $a^{-1} = 1/a$
    \item $(a^m)^n = a^{mn}$
    \item $a^m a^n = a^{m+n}$
    \item $e^x \ge 1 + x$
\end{itemize}
\vspace{8pt}
\textbf{Logarithms}
\begin{itemize}
    \item \textbf{Binary log:} $\log n = \log_2 n$
    \item \textbf{Natural log:} $\ln n = \log_e n$
    \item \textbf{Exponentiation:} $\log^k n = (\log n)^k$
    \item \textbf{Composition:} $\log \log n = \log(\log n)$
    \item $a = b^{\log_b a}$
    \item $\log_c (ab) = \log_c a + \log_c b$
    \item $\log_b a^c = c \log_b a$
    \item \textbf{Change of Base Formula:} $\log_b a = \frac{\log_c a}{\log_c b}$
    \item $\log_b (1/a) = -\log_b a$
    \item $\log_b a = \frac{1}{\log_a b}$
    \item $a^{\log_b c} = c^{\log_b a}$
\end{itemize}
\vspace{8pt}
\textbf{Stirling's approximation}
\begin{itemize}
    \item Estimate for the factorial function
    \item $n! = \sqrt{2\pi n} \left(\frac{n}{e}\right)^n (1 + O(1/n))$
    \item $\log(n!) \in \Theta(n \log n)$
\end{itemize}
\vspace{8pt}
\textbf{Summation and series}
\begin{itemize}
    \item Finite GP: $\sum_{k=0}^{n} x^k = 1 + x + x^2 + \cdots + x^n = \frac{x^{n+1} - 1}{x - 1}$
    \item Infinite GP: $\sum_{k=0}^{\infty} x^k = \frac{1}{1 - x}$
    \item AP: $\sum_{k=0}^{n} k = 1 + 2 + \cdots + n = \frac{n(n + 1)}{2}$, $\Theta(n^2)$
    \item Finite AP term: $T_k = a + (k-1)d$ 
    \item Finite AP sum: $S_n = \frac{n}{2}(2a + (n - 1)d)$
    \item Finite GP term: $T_k = ar^k$
    \item Finite GP sum: $S_n = \frac{a(1 -  r^n)}{1-r}$
    \item Harmonic series: $\sum_{i=1}^{n} \frac{1}{i} = \ln(n)+\gamma + o(1) \in \theta(\log n)$
    \item Reciprocal Log: $\sum_{k=2}^{n} \frac{1}{\log k} = \theta(n / \log n)$ 
    \item Harmonic over log: $\sum_{k=2}^{n} \frac{1}{k \log k} = \theta(\log \log n)$
    \item Squared log: $\sum_{k=2}^{n} \frac{1}{k(\log k)^p}$ 
    \begin{itemize}
        \item $p > 1$: $\theta (1)$
        \item $p = 1$: $\theta (\log \log n)$
        \item $0 < p < 1$: $\theta((\log n)^{1-p})$
    \end{itemize}
\end{itemize}
\vspace{8pt}
\textbf{Generalised Harmonic Series} \\
$\sum_{i=1}^{n} \frac{1}{i^p}$
\begin{itemize}
    \item $p > 1$: $\theta (1)$
    \item $p = 1$: $\theta (\log n)$
    \item $0 < p < 1$: $\theta(n^{1-p})$
\end{itemize}
\vspace{8pt}
\textbf{L'Hospital's Rule}
\begin{itemize}
    \item technique for evaluating indeterminate limits
    \item if $\lim_{x\to\infty} f(x) = \infty$ and $\lim_{x\to\infty} g(x) = \infty$, then:
    \begin{itemize}
        \item $\lim_{x\to\infty} \frac{f(x)}{g(x)} = \lim_{x\to\infty} \frac{f'(x)}{g'(x)}$
    \end{itemize}
    \item where $f'(x)$ and $g'(x)$ are the derivatives of $f(x)$ and $g(x)$.
    \item $n \ln n \in o(n^2)$, meaning $n \ln n$ grows strictly slower than $n^2$
\end{itemize}
\vspace{8pt}
\textbf{Useful inequalities}
\begin{itemize}
    \item $1 + x \leq e^x$
\end{itemize}
\vspace{8pt}
\textbf{Sample Space and Event}
\begin{itemize}
    \item \textbf{Sample Space ($S$):} The set of all possible outcomes (elementary events) of an experiment.
    \item \textbf{Event ($A$):} Any subset of the sample space.
    \item \textbf{Complement of an Event ($\bar{A}$):} The set of all outcomes in $S$ that are not in $A$, defined as $\bar{A} = S - A$.
\end{itemize}
\vspace{8pt}
\textbf{Probability Distribution} \\
A probability distribution $\Pr()$ is a mapping from events to real numbers that satisfies three axioms:
\begin{enumerate}
    \item \textbf{Non-negativity:} $\Pr(A) \ge 0$ for all events $A$.
    \item \textbf{Normalization:} $\Pr(S) = 1$.
    \item \textbf{Additivity:} $\Pr(A \cup B) = \Pr(A) + \Pr(B)$ for mutually exclusive (disjoint) events $A$ and $B$.
\end{enumerate}
\vspace{8pt}
\textbf{General Probability Rules}
\begin{itemize}
    \item \textbf{Inclusion-Exclusion Principle (for two events):} $\Pr(A \cup B) = \Pr(A) + \Pr(B) - \Pr(A \cap B)$.
    \item \textbf{Independence:} Two events $A$ and $B$ are independent if and only if $\Pr(A \cap B) = \Pr(A) \cdot \Pr(B)$.
\end{itemize}
\vspace{8pt}
\textbf{Conditional Probability} \\
The conditional probability of event $A$ given event $B$ (where $\Pr(B) \ne 0$) is defined as:
\begin{itemize}
    \item $\Pr(A|B) = \frac{\Pr(A \cap B)}{\Pr(B)}$
\end{itemize}
\vspace{8pt}
\textbf{Bayes' Theorem} \\
$\Pr(A|B) = \frac{\Pr(A) \Pr(B|A)}{\Pr(B)}$.
\vspace{8pt}
The denominator $\Pr(B)$ can be expanded using the Law of Total Probability (for an event $A$ and its complement $\bar{A}$):
\[ \Pr(A|B) = \frac{\Pr(A) \cdot \Pr(B|A)}{\Pr(A) \cdot \Pr(B|A) + \Pr(\bar{A}) \cdot \Pr(B|\bar{A})} \]
\vspace{8pt}
\textbf{Random Variable} \\
A function that maps the outcomes in the sample space ($S$) to real numbers.
\begin{itemize}
    \item The function $f(x) = \Pr(X = x)$ is the \textbf{probability density function} (or probability mass function for discrete variables) of $X$.
\end{itemize}
\vspace{8pt}
\textbf{Expectation (Mean)} \\
The expectation or mean of a random variable $X$ is the weighted average of all possible values of $X$:
\begin{itemize}
    \item $E(X) = \sum_{x} x \cdot \Pr(X = x)$
\end{itemize}
\vspace{8pt}
\textbf{Linearity of Expectation} \\
Expectation is a linear operator:
\begin{itemize}
    \item $E(X + Y) = E(X) + E(Y)$
    \item $E(aX) = a E(X)$
    \item \textbf{For independent variables:} $E(XY) = E(X) \cdot E(Y)$
\end{itemize}
\vspace{8pt}
\textbf{Event Operations}
\begin{enumerate}
    \item $A \cap A' = \emptyset$
    \item $A \cap \emptyset = \emptyset$
    \item $A \cup A' = S$
    \item $(A')' = A$
    \item $A \cup (B \cap C) = (A \cup B) \cap (A \cup C)$
    \item $A \cap (B \cup C) = (A \cap B) \cup (A \cap C)$
    \item $A \cup B = A \cup (B \cap A')$
    \item $A = (A \cap B) \cup (A \cap B')$
    \item $(A_1 \cup A_2 \cup \dots \cup A_n)' = A_1' \cap A_2' \cap \dots \cap A_n'$
    \item $(A \cup B)' = A' \cap B'$
    \item $(A_1 \cap A_2 \cap \dots \cap A_n)' = A_1' \cup A_2' \cup \dots \cup A_n'$
    \item $(A \cap B)' = A' \cup B'$
\end{enumerate}
\vspace{8pt}
\textbf{Permutations and Combinations} \\
\begin{enumerate}
    \item $P^n_r=\frac{n!}{(n-r)!}=n(n-1)(n-2)\dots (n-(r-1))$ \\
    \item $n \choose {r}$ = $\frac{n!}{r!(n-r)!}$ \\
    \item $n\choose r$ = $n \choose (n-r)$
\end{enumerate}
\vspace{8pt}
\textbf{Probability}
\begin{itemize}
    \item $P(A)=P(A\cap{B})+P(A\cap \overline{B})$  
    \item $P(A_1 \cap A_2 \cap \dots \cap A_n) = P(A_1) + P(A_2) + \dots + P(A_n)$, if $A_i \cap A_j = \emptyset$ for any $i \neq j$
    \item $P(A \cup B) = P(A) + P(B) - P(A \cap B)$ ------- Principle of Inclusion and Exclusion
    \item If $A \subset B$, then $P(A) \leq P(B)$ 
    \item $A$ and $B$ are independent $\iff$ $P(A \cap B) = P(A)P(B)$ ------- $i.e$ $A\perp B$   
    \item Suppose P(A) > 0, P(B) > 0. If A and B are independent, then A and B are not mutually exclusive.
    \item Suppose P(A) > 0, P(B) > 0. If A and B are mutually exclusive, then A and B are dependent. 
    \item $S$ and $\emptyset$ are independent of any other event.
    \item If A $\perp$ B, then A $\perp$ B', A' $\perp$ B, and A' $\perp$ B'.
    \item $A \perp B \iff P(B | A) = \frac{P(A \cap B)}{P(A)} = \frac{P(A)P(B)}{P(A)} = P(B)$, if $P(A) \neq 0$
    \item Contained: $A\subset{B}$ 
    \item Equivalent: $A=B$
    \item Mutually Exclusive: $A\cap{B}=\emptyset$ 
    \item Law of Total Probability: $P(B) = \sum^n_{i=1}P(B\cap A_i)=\sum^n_{i=1}P(A_i)P(B|A_i)$ 
    \item $P(A|B)=\frac{P({A}\cap{B})}{P(B)}$ 
    \item $P(A|B) = \frac{P(B|A)P(A)}{P(B)}$
    \item $P(A_k|B)=\frac{P(A_k)P(B|A_k)}{\sum^n_{i=1}P(A_i)P(B|A_k)}$
    \item $P(A|B)=1-P(\overline{A}|B)$ 
    \item $P(A)=P(A|B)P(B)+P(A|\overline{B})P(\overline{B})$ 
\end{itemize}
\vspace{8pt}
\textbf{Bernoulli Trials}
\begin{itemize}
    \item A \textbf{Bernoulli Trial} is an experiment with exactly two outcomes: success (with probability $p$) and failure (with probability $q = 1 - p$).
\end{itemize}
\vspace{8pt}
\textbf{Geometric Distribution} \\
If $X$ is the number of independent Bernoulli trials needed to obtain the first success:
\begin{itemize}
    \item \textbf{Probability:} $\Pr(X = k) = q^{k-1} p$
    \item \textbf{Expectation:} $E(X) = \frac{1}{p}$
\end{itemize}
\vspace{8pt}
\textbf{Binomial Distribution} \\
If $X$ is the number of successes in $n$ independent Bernoulli trials:
\begin{itemize}
    \item \textbf{Probability:} $\Pr(X = k) = \binom{n}{k} p^k q^{n-k}$
\end{itemize}

\hrulefill

\section*{1. Asymptotic Analysis}

\textbf{Big-O (g is upper bound of f)}: $f \in O(g)$
\begin{itemize}
    \item Definition: $\exists c > 0 \land n_0 > 0 \space \text{s.t} \space \forall n \geq n_0: 0 \leq f(n) \leq c \cdot g(n)$
    \item Theorem+Proof: 
    \begin{enumerate}
        \item $\lim_{n\rightarrow \infty} \frac{f(n)}{g(n)} < \infty$
        \item $\lim_{n\rightarrow \infty} \frac{f(n)}{g(n)} = L, L \in \mathbb{R}$
        \item $\forall \epsilon > 0, \exists n_0 \space \text{s.t} \space \forall n \geq n_0, \frac{f(n)}{g(n)} - L < \epsilon$
    \end{enumerate}
\end{itemize}
\vspace{8pt}
\textbf{$\Omega$ (g is lower bound of f):} $f \in \Omega (g(n))$
\begin{itemize}
    \item Definition: $\exists c > 0 \land n_0 > 0 \space \text{s.t} \space \forall n \geq n_0: 0 \leq c \cdot g(n) \leq f(n)$ 
    \item Theorem+Proof:
    \begin{enumerate}
        \item $\lim_{n\rightarrow \infty} \frac{f(n)}{g(n)} > 0$
        \item $\lim_{n\rightarrow \infty} \frac{f(n)}{g(n)} = L, L > 0$
        \item $\lim_{n\rightarrow \infty} \frac{g(n)}{f(n)} = \frac{1}{L}$
        \item $\lim_{n\rightarrow \infty} \frac{f(n)}{g(n)} < \infty$
    \end{enumerate}
\end{itemize}
\vspace{8pt}
$\Theta$ \textbf{(g is tight bound of f)}: $f \in \Theta (g(n))$
\begin{itemize}
    \item Definition: $\exists c1, c2 > 0 \land n_0 > 0 \space \text{s.t} \space \forall n \geq n_0: 0 \leq c_1 \cdot g(n) \leq f(n) \leq c2 \cdot g(n)$
    \item Theorem+Proof:
    \begin{enumerate}
        \item $0 < \lim_{n\rightarrow \infty} \frac{f(n)}{g(n)} < \infty$
        \item $\Theta = O \cap \Omega$
    \end{enumerate}
\end{itemize}
\vspace{8pt}
\textbf{Little-o (g is strict upper bound of f)}: $f \in o(g(n))$
\begin{itemize}
    \item Definition: $\forall c > 0, \exists n_0 > 0 \space \text{s.t} \space \forall n \geq n_0: 0 \leq f(n)  < c \cdot g(n)$
    \item Theorem+Proof: 
    \begin{enumerate}
        \item $\lim_{n\rightarrow \infty} \frac{f(n)}{g(n)} = 0$
        \item $\forall \epsilon > 0,  \exists n_0 > 0 \space \text{s.t} \space \forall n \geq n_0: \frac{f(n)}{g(n)} < \epsilon$
    \end{enumerate}
\end{itemize}
\vspace{8pt}
$\omega$ \textbf{(g is strict lower bound of f)}: $f \in \omega(g(n))$
\begin{itemize}
    \item Definition: $\forall c > 0, \exists n_0 > 0 \space \text{s.t} \space \forall n \geq n_0: 0 \leq c \cdot g(n) < f(n)$
    \item Theorem+Proof:
    \begin{enumerate}
        \item $\lim_{n\rightarrow \infty} \frac{f(n)}{g(n)} = \infty$
        \item $\forall M > 0, \exists n_0 \space \text{s.t} \space \forall n \geq n_0: \frac{f(n)}{g(n)} > M$
    \end{enumerate}
\end{itemize}
\vspace{8pt}
\textbf{Properties} \\
\textbf{1. Reflexivity}
\begin{itemize}
    \item $f(n) \in O(f(n))$ (Choose c = 1, $\forall n, f(n) \leq 1 \cdot f(n)$)
    \item $f(n) \in \Omega(f(n))$ (Choose c = 1, $\forall n, f(n) \geq 1 \cdot f(n)$)
    \item $f(n) \in \Theta(f(n))$ (Chose c1 = c2 = 1)
\end{itemize}
\vspace{8pt}
\textbf{2. Transitivity}
\begin{itemize}
    \item $f(n) \in O(g(n)) \land g(n) \in O(h(n)) \implies f(n) \in O(h(n))$
    \item $f(n) \in \Omega(g(n)) \land g(n) \in \Omega(h(n)) \implies f(n) \in \Omega(h(n))$
    \item $f(n) \in \Theta(g(n)) \land g(n) \in \Theta(h(n)) \implies f(n) \in \Theta(h(n))$
    \item $f(n) \in o(g(n)) \land g(n) \in o(h(n)) \implies f(n) \in o(h(n))$
    \item $f(n) \in \omega(g(n)) \land g(n) \in \omega(h(n)) \implies f(n) \in \omega(h(n))$
\end{itemize}
\vspace{8pt}
\textbf{3. Symmetry}
\begin{itemize}
    \item $f(n) \in \Theta(g(n)) \iff g(n) \in \Theta(f(n))$
\end{itemize}
\vspace{8pt}
\textbf{4. Complementarity}
\begin{itemize}
    \item $f(n) \in O(g(n)) \iff g(n) \in \Omega(f(n))$
    \item $f(n) \in o(g(n)) \iff g(n) \in \omega(f(n))$
\end{itemize}

\hrulefill

\section*{2. Recurrences and Master Theorem}

\textbf{Remarks}:
\begin{itemize}
    \item If $f(n) \in O(n)$, then $\exists c > 0, n_0 > 0$ s.t. $f(n) \leq cn$ if $n \geq n_0$. For \textbf{upper bound}, replace $\Theta(n)$ with $cn$.
    \item If $f(n) \in \Omega(n)$, then $\exists c > 0, n_0 > 0$ s.t. $f(n) \geq cn$ if $n \geq n_0$. For \textbf{lower bound}, replace $\Theta(n)$ with $cn$.
\end{itemize}
\vspace{8pt}
\textbf{(1) Telescoping Series}
\begin{itemize}
    \item Find a way to cancel out each other
    \item e.g $T(n)/n = \frac{2}{n}T(\frac{n}{2}) + 1 = T(\frac{n}{2})/\frac{n}{2} + 1$ 
\end{itemize}
\vspace{8pt}
\textbf{(2) Substitution Method}
\begin{itemize}
    \item Guess a solution, then verify by induction
\end{itemize}
\vspace{8pt}
\textbf{(3) Recursion Tree}
\begin{itemize}
    \item Find how many levels and leaves
\end{itemize}
\vspace{8pt}
\textbf{(4) Master Theorem}
\begin{itemize}
    \item Generic form: $T(n) = aT(n / b) + f(n), a \geq 1, b \geq 1$
    \item Classify work: $f(n)$ (root) vs $n^d$ (leaves) where $d = \log_b a$ is the \textbf{critical exponent}.
    \item \textbf{Case 1}: $f(n) \in O(n^{d-\epsilon})$ for some $\epsilon > 0$. Ans: $T(n) \in \Theta(n^d)$.
    \item \textbf{Case 2}: $f(n) \in \Theta(n^d \log^k n)$ for some $k \geq 0$. Ans: $T(n) \in \Theta(n^d \log^{k+1} n)$.
    \item \textbf{Case 3}: $f(n) \in \Omega(n^{d+\epsilon})$ for some $\epsilon > 0$ and regularity condition $af(n/b) \leq cf(n)$ for $c < 1$. Ans: $T(n) \in \Theta(f(n))$.
\end{itemize}

\hrulefill

\section*{3. Proof of correctness}

\textbf{Iterative algorithm}
\begin{enumerate}
    \item Loop invariant: Induction hypothesis
    \item Initialisation: Base case
    \item Maintenance: Inductive step
    \item Termination: Proof by induction implies the correctness
\end{enumerate}
\vspace{8pt}
\textbf{Recursive Algorithms}
\begin{enumerate}
    \item Base case
    \item Inductive step: Assume correct for size $\leq n$, show for $n$.
\end{enumerate}
\vspace{8pt}
\textbf{Fibonacci}
\begin{enumerate}
    \item Loop invariant: At start of iteration $i$: $prev2 = Fib(i - 2), prev1 = Fib(i - 1)$
    \item Initialisation: $i = 2: prev2 = Fib(0) = 0, prev1 = Fib(1) = 1$
    \item Maintenance: If true at $i$, then at end $prev2 = Fib(i - 1), prev1 = Fib(i)$
    \item Termination: Returns $Fib(n)$
\end{enumerate}
\vspace{8pt}
\textbf{Selection sort}
\begin{enumerate}
    \item Loop invariant: $A[1..j-1]$ is sorted and ($x \leq y$ for all $x \in A[1..j-1]$ and $y \in A[j..n]$)
    \item Initialisation: True because nothing in LHS
    \item Maintenance: Take new element from $A[j..n]$ and swap, invariant maintained
    \item Termination: New element inserted is bigger than that already inserted, thus sorted.
\end{enumerate}
\vspace{8pt}
\textbf{Dijkstra} \\
\textbf{Definitions}
\begin{enumerate}
    \item $s$: Source vertex
    \item $R$: Set of already correct computed distances
    \item $dist(u, v)$: Distance from $u$ to $v$
    \item $w(u, v)$: Edge cost from $u$ to $v$
    \item $V$: Set of all vertices
    \item $D_v$: Set of all possible path distances from $s$ to $v$ 
        \begin{itemize}
            \item i.e $D_v = \{\forall u \in R : d(u) + w(u, v)\}$
        \end{itemize}
    \item $d(v)$: $min(D_v)$
\end{enumerate}
\vspace{8pt}
\textbf{Algorithm}
\begin{enumerate}
    \item Base case: $dist(s,s) = 0$, $R = \{s\}$
    \item Loop: While $R \neq V$:
        \begin{enumerate}
            \item Select $v \in V$ such that $v$ is not in $R$ and $v$ minimises $D_v$
            \item Update $d(v) = min(D_v)$
            \item Add $v$ to $R$
        \end{enumerate}
\end{enumerate}
\vspace{8pt}
\textbf{Loop invariant}: At the start of every iteration, $\forall u \in R, dist(s, v) = min(D_v)$
\vspace{8pt} \\
\textbf{Proof of correctness of the loop invariant}
\begin{enumerate}
    \item Initialisation: $dist(s, s) = d(s) = 0$
    \item Termination: 
        \begin{enumerate}
            \item At the termination point, in each iteration of the while loop, we add one extra vertex to R
            \item Every vertex $u \in R = d(u)$, and $R = V$ at the end of the computation (i.e every vertex has the shortest distance from $s$ computed correctly)
            \item So at the end of the algorithm, all the vertices will be in $R$ and $R$ represents the shortest path from $s$ to every other vertices at termination
        \end{enumerate}
    \item Maintenance (i.e how each loop iteration maintains the invariant)
        \begin{enumerate}
            \item At the start of the loop iteration, the algorithm chooses a new vertex $v$ to add to $R$ such that $min(D_v)$ is minimised (step 2.1 in the algorithm) (now to prove that this maintains the invariant)
            \item Consider the shortest path $P$ from $s$ to $v$. 
                \begin{itemize}
                    \item Visual: $s \rightarrow \dots \rightarrow y \rightarrow v$
                \end{itemize}
            \item Claim: Every vertex in $P$ is in $R$ except for $v$ 
                \begin{itemize}
                    \item Visual: $(s \rightarrow \dots \rightarrow y) \rightarrow v$
                    \item Need to prove that this claim is true, because if this claim is true, then it follows that the loop invariant is maintained (i.e $dist(s,v) = d(v)$)
                \end{itemize}
            \item Suppose the claim (3.3) is false:
                \begin{enumerate}
                    \item Then there exists a vertex $x$ before $v$ and immediately after $y$
                        \begin{itemize}
                            \item Visual: $(s \rightarrow \dots \rightarrow y) \rightarrow x \rightarrow \dots \rightarrow v$
                        \end{itemize}
                    \item The true shortest path from $s$ to $v$ = $dist(s, v)$ is in $D_v$. So $min(D_v) \geq dist(s,v)$
                    \item From the visual 3.4.1.1, we have $dist(s, v) > dist(s, y) + w(y, x)$
                    \item Since $y$ is in $R$, $dist(s,y) = min(D_y)$ 
                    \item $dist(s, y) + w(y, x) = min(D_y) + w(y, x) \geq min(D_x)$
                    \item Full inequality: $min(D_v) \geq dist(s,v) > dist(s,y) + w(y,x) \geq min(D_x)$ 
                    \item Which implies $min(D_v) \geq min(D_x)$
                    \item So the algorithm would've selected $x$ instead of $v$. This contradicts (3.1).
                \end{enumerate}
            \item Thus, the claim (3.3) must be true, which implies:
                \begin{enumerate}
                    \item $min(D_v) \geq dist(s, v) = d(y) + w(y, v) \geq min(D_v)$
                    \item $dist(d, v) = min(D_v)$
                    \item Visual: $(s \rightarrow \dots \rightarrow y) \rightarrow v$
                \end{enumerate}
            \item Thus, the loop invariant is satisfied ($\forall u \in R, dist(s, u) = min(D_u)$)
        \end{enumerate}
\end{enumerate}
\vspace{8pt}
\textbf{Efficiency of Djikstra}
\begin{itemize}
    \item Priority queue implementation.
\end{itemize}

\begin{tabular}{|l|l|l|l|l|}
\hline
Data Structure & Insert & Extr-min & Decr-key & Time complexity \\ \hline
Binary Heap & $O(\log n)$ & $O(\log n)$ & $O(\log n)$ & $O((|E| + |V|) \log |V|)$ \\ \hline
Fibonacci Heap & $O(1)$ & $O(\log n)$ & $O(1)$ & $O(|E| + |V|\log |V|)$ \\ \hline
\end{tabular}
\vspace{8pt} \\
\textbf{Searching in a sorted array} \\
\vspace{8pt}

\hrulefill

\section*{4. Sorting}


\hrulefill

\section*{5. Randomized Algorithms}


\end{multicols*}

\end{document}
